\documentclass{report}
\usepackage{amsfonts, amsmath}
\newcommand\R{{\mathbb R}}
\newcommand\Q{{\mathbb Q}}
\newcommand\N{{\mathbb N}}
\newcommand\Z{{\mathbb Z}}
\newcommand\Or{{\mathcal O}}
\newcommand\ve{\varepsilon}
%\pagestyle{empty}
\begin{document}
\large
\centerline{\Huge Fisica Matematica II}
\renewcommand{\thefootnote}{ } 
\renewcommand{\thefootnote}{\arabic{footnote}} 
\vskip.1cm
\centerline{\Large Soluzioni appello 10-07-2003}
\vskip1cm
\begin{enumerate}
\item Si controlla immediatamente che la condizione iniziale non \`e
una caratteristica. Ci si riconduce dunque al sistema
\[
\begin{split}
&\dot x=y^2-z^2\\
&\dot y=-xy\\
&\dot z=xz\\
&z(0)=s\\
&y(0)=s\\
&x(0)=s.
\end{split}
\]
Si ottiene:
\[
\dot y y +\dot z z=x(z^2-y^2)=-\dot xx;\;\; \dot yxz=-\dot zxy.
\]
Da cui si ahnno i due integrali del moto
\[
x^2+y^2+z^2=3s^2;\;\; yz=s^2
\]
cio\`e
\[
u=\frac{3y-\sqrt{5y^2-4x^2}}{2}.
\]
\item Si calcola
\[
\int_Q|\nabla u|^2=\int_Q\text{div }u\nabla u-u\Delta u
=\int_{\partial Q}u\frac {\partial u}{\partial n} +\int _Q\lambda u^2=
\lambda\int _Q u^2
\]
Dunque $\lambda >0$. Nota inoltre che $u=\sin\pi x\;\sin \pi y$ risolve
il probelma e quindi si puo avere $\lambda =2\pi^2$. In effetti questo
il valore minimo possibile di $\lambda$. Pe verificarlo si noti che
una soluzione deve essere infinitamente differenziabile e quindi si puo
scrivere come una serie seno di Fourier
\[
u(x,y)=\sum_{n,m\geq 1}u_{nm}\sin n\pi x \sin m \pi y.
\] 
Sostituendo nell'equazione si ottine che esistono soluzioni solo se
$\lambda=2(n^2+m^2)\pi^2$. 
\item La soluzione \`e data da
\[
u_a(x,t)=\frac 1{\sqrt{4\pi ta}}\int_{\R}
e^{-\frac{(x-y)^2}{4t}-\frac{x^2}a}dy.
\]
da cui $\lim_{a\to\infty}u_a(x,1)=0$.
\item Si cerca una sol del tipo
\[
u(x,t)=\sum_{n=1}^\infty c_n(t)\sin n\pi x,
\]
che soddisfa le condizioni al contorno.
Sostituendo si ottiene
\[
\begin{split}
&c_n''(t)=-n^2c_n(t)+b_ne^{-t}\\
&c_n(0)=0;\;\; c'_n(0)=0,
\end{split}
\]
where $1=\sum_{n=1}^\infty b_n\sin n\pi x$. Cio\`e
\[
b_n=2\int_0^1x\sin n\pi x \, dx =\frac 2n [1-(-1)^n].
\]
La soluzione \`e dunque $c_n=0$ per $n$ pari e 
\[
c_n(t)=\frac{b_n}n\int_0^te^{-s}\sin n(t-s)\;ds.
\]
Finalmente,
\[
\frac d{dt} E(t)=e^{-t}\int_0^1u_t(x,t)dx.
\]
Dunque il limite richiesto \`e nullo.
\end{enumerate}
\end{document}
